\subsection{Методы структурной оптимизации распределения. Metis}
METIS (от MEtis Tournament Inspired Strategy) \cite{MetisOG} - это набор программных библиотек и утилит, разработанных в Университете Миннесоты, для разбиения больших неориентированных графов и разрежения матриц. На данный момент считается эталонным решением структурной оптимизации распределения.

Ключевая идея METIS - многоуровневая парадигма (Multilevel Paradigm). Вместо того чтобы работать с огромным исходным графом напрямую, METIS последовательно его упрощает, находит разбиение для маленького графа и затем интерполирует это разбиение обратно на исходный граф. В силу этой многоуровневой парадигмы в контексте Metis будет удобно рассмотреть различные методы, которые могут применяться на каждом уровне.

Алгоритм состоит из трёх этапов: 
\begin{enumerate}
    \item Схлопывание/Упрощение/Свёртка (англ. Coarsening) - уменьшение графа схлопыванием вершин в граф $G_l$. \\
    Исходный граф $G_0$ последовательно преобразуется во всё более мелкие графы $G_1, G_2, ..., G_l$. На каждом шаге пары смежных вершин "схлопываются/сворачиваются" в одну супервершину.
    \item Распределение (англ. Partitioning) - распределение графа $G_l$. Поскольку граф $G_l$ очень мал, для его разбиения можно использовать даже очень дорогие (вычислительно сложные) алгоритмы.
    \item Развёртка (англ. Uncoarsening) - развёртка графа $G_l$, то есть действие обратное первому этапу. Также происходит улучшение (англ. refinement) - оптимизация разбиения в процессе развёртки. Порядок следующий:
    \begin{enumerate}
        \item Разбиение, найденное для $G_l$, проецируется на граф $G_{m-1}$. Так как каждая вершина в $G_{m-1}$ соответствовала вершине в $G_l$, разбиение определяется автоматически.
        \item Улучшение, то есть переброс вершин между шардами для получения лучшего распределения.
        \item Первые два пункта повторяются пока не будет получено разбиение для $G_0$.
    \end{enumerate} Этап хорошо подвергается параллелизации.
    
\end{enumerate}

\subsubsection{Алгоритмы свёртки}

Свёртка в Metis сводится к поиску максимального паросочетания, но не наибольшего, так как это слишком затратное действие. и дальнейшего схлопывания рёбер паросочетания.

В базовом методе рассматриваются 4 алгоритма поиска максимального паросочетания:
\begin{enumerate}
    \item Случайное (англ. Random Matching - RM) - берётся случайное ребро, затем случайное ребро, несмежное с ним, после чего ребро не связанное с выбранными ранее, и так далее пока таких рёбер не останется. Метод прост и эффективен для локальной задачи, но для уменьшения общей мощности разрезов в дальнейшем подходит плохо.
    \item Паросочетание тяжёлых рёбер (англ. Heavy Edge Matching - HEM) - выбор рёбер максимального веса по очереди. Экспериментально было установлено, что это приводит к хорошему уменьшению общей мощности разрезов. В Metis этот алгоритм используется как основной.
    \item Паросочетание лёгких рёбер (англ. Light Edge Matching - LEM) - выбор рёбер минимального веса по очереди. Алгоритм приводит к большему общему весу рёбер у свёрнутого графа, что может положительно влиять на качество улучшения некоторыми алгоритмами.
    \item Паросочетание тяжёлых клик (англ. Heavy Clique Matching - HCM) - объединяет вершины, максимизируя плотность ребер создаваемого подграфа.

В отличие от HEM, который выбирает тяжелые ребра, HCM оценивает, насколько две вершины (точнее вершины, которые свернулись в них) близки к формированию клики. Для вершин $u$ и $v$ вычисляется плотность рёбер (англ. edge density) по формуле:
    \begin{equation*}
        \text{density}(u,v) = \frac{2(ce(u) + ce(v) + ew(u,v))}{(vw(u) + vw(v))(vw(u) + vw(v) - 1)}
    \end{equation*}, где $ce(u)$, $ce(v)$ - веса уже схлопнутых рёбер в вершинах, $ew(u,v)$ - вес ребра между $u$ и $v$, $vw(u)$,$ vw(v)$ - сумма степеней вершин, схлопнутых в данные.
\end{enumerate}

Алгоритм схож с HEM, но учитывает и вершины изначального графа.

\subsubsection{Алгоритмы распределения}
Этап распределения является решением поставленной выше задачи для свёрнутого графа. В Metis рекурсивно используются алгоритмы бисекции графа, поэтому число шардов $k$ должно быть степенью двойки, другоие $k$ приведут к плохой балансировке.

На этом этапе также рассматриваются 4 алгоритма:
\begin{enumerate}
    \item Спектральная бисекция (англ. Spectral Bisection - SB) - Вычисляет собственный вектор Фидлера $y$ матрицы Лапласиана $Q=D-AQ=D-A$, где:
    \begin{equation*}
        a_{ij} = \begin{cases} 
ew(v_i, v_j) & \text{если } (v_i, v_j) \in E_m \\
0 & \text{иначе}
\end{cases}
    \end{equation*}
    \begin{equation*}
    d_{ii} = \sum_{(v_i, v_j) \in E_m} ew(v_i, v_j)
    \end{equation*}

    Разбиение: $P[j]=1$ если $y_j\le r$, и $P[j]=2$ в противном случае, где $r$ - выбранная взвешенная медиана $y_j$.
    \item Алгоритм Кернигана-Лина (англ. Kernighan-Lin - KL) - алгоритм, работающий на идее улучшения разбиения путём обмена вершин между шардами. Для этого определена следующая функция улучшения распределения:
    \begin{equation*}
        g_v = \sum_{(v,u) \in E \wedge P[v] \neq P[u]} w(v,u) - \sum_{(v,u) \in E \wedge P[v] = P[u]} w(v,u)
    \end{equation*}
    при положительном $g_v$ перемещение вершины $v$ приведёт к улучшению распределения. За одну итерацию алгоритм проходится по всем вершинам и пытается переместить все, итерации прерываются как только никакое перемещение не приведёт к улучшению, либо когда будет достигнуто предельное число операций. Эксперименты показали, что для достаточно хорошего распределения как правило хватает 5-10 итераций.\cite{MetisOG}. Для избегания работы "вхолостую" итерация прерывается при 50 перемещений без улучшения, что хорошо показало себя при экспериментах. Также стоит отметить, что ещё можно ограничить алгоритм на минимальное кол-во перемещений за итераций. 
    
    Функция улучшения может быть рассчитана для всех вершин в начале итерации и при перемещении вершины обновляться только для соседей перемещённой вершины для оптимизации. 
    \item Алгоритм роста графа (англ. Graph Growing - GGP) - выбирается случайная вершина и от неё как при поиске в ширину обходятся вершины, формируя подграф, пока не будет набрана ровно половина вершин. Этот подграф будет первым шардом, остальные вторым. Далее рекомендуется применить KL для улучшения, что не займёт много времени в силу малого размера свёрнутого графа. Также стоит отметить, что вполне можно делить не пополам, а на другие равные части и применять KL к шардам попарно.
    \item Жадный алгоритм роста графа (англ. Greedy Graph Growing - GGGP) - так как для предыдущего алгоритма и так рекомендуется применять KL после бисекции можно сразу при обходе графа идти по вершинам, которые будут давать наибольшее улучшение.
\end{enumerate}

\subsubsection{Алгоритмы уточнения}
В Metis предлагаются два схожих алгоритма:
\begin{enumerate}
    \item  Алгоритм уточнения Кернигана-Лина  - предлагается продолжать использовать KL на каждом новом развёрнутом графе $G_i$. Так как $G_{i+1}$ Уже был хорошо распределён большого кол-ва перемещений не потребуется. Эксперименты показали, что на первой итерации получается наибольшее улучшение, поэтому применение только одной итерации выделяют как алгоритм KL(1) как значительном менее затратный, чем KL.
    \item Граничный алгоритм уточнения Кернигана-Лина (англ. Boundary Kernighan-Lin Refinement - BKL) - так как итерации прерываются большинство вычислений тратятся впустую, особенно в KL(1), поэтому предлагается прменять KL только к "граничным" вершинам, то есть к тем, которые смежны с вершинами из другого шарда. Это позволяет сильно уменьшить вычислительные затраты, что в свою очередь позволяет чаще делать несколько итераций, что ведёт к идее алгоритма BKL(*,1) - применять BKL пока вершин немного и перейти на BKL(1), когда их станет слишком много.
\end{enumerate}

Из всего вышесказанного стоит выделить алгоритм уточнения Кернигана-Лина, и его модификацию -- граничный алгоритм уточнения Кернигана-Лина --, позволяющие улучшать уже существующее распределение, если оно было каким-либо образом модифицировано. Больше о его возможном применении дальше.